\chapter{Complexity — Logspace Uniform Example}
%%% generated by the blueprint pipeline from TCSlib.Complexity.CircuitComplexity.LogspaceUniformExample
\section{Overview}

Declarations of the Lean module \texttt{TCSlib.Complexity.CircuitComplexity.LogspaceUniformExample}.

\section{Declarations}

\begin{definition}[Adjacent-value language]
\lean{BoolCircuit.AdjExample.leLang}
\label{BoolCircuit.AdjExample.leLang}
\leanok
\uses{Turing.pairEncode}
For Boolean values $a$ and $b$, let $L_{a,b}$ consist of the strings encoding a pair
whose first component is a string of $n$ ones and whose second component is the binary
representation of some integer $i$, for some $n,i\in\mathbb{N}$, such that either $i<n$
and $a$ is true, or $i=n$ and $b$ is true.
\end{definition}

\begin{definition}[Comparison machine states]
\lean{BoolCircuit.AdjExample.Lb}
\label{BoolCircuit.AdjExample.Lb}
\leanok
Let the state set consist of eight distinct states: start, test, compare, increment,
equality-check, return-$a$, return-$b$, and return-false.
\end{definition}

\begin{definition}[Adjacent-value register program]
\lean{BoolCircuit.AdjExample.leA}
\label{BoolCircuit.AdjExample.leA}
\leanok
\uses{BoolCircuit.AdjExample.Lb, Complexity.LogProg.ARM}
For Boolean values $a$ and $b$, define a one-register program with one available decider
and the eight states start, test, compare, increment, equality-check, return-$a$,
return-$b$, and return-false. At start it validates the input as a unary prefix followed
by a canonical binary payload; at test it calls the decider on the unary-prefix mode
with register $0$ as its argument, branching to compare on acceptance and equality-check
otherwise. At compare it tests whether the input payload equals the binary
representation of register $0$, returning $a$ if so and otherwise incrementing the
register and returning to test; at equality-check it makes the same payload comparison,
returning $b$ if equal and false otherwise. The three return states halt with outputs
$a$, $b$, and false, respectively.
\end{definition}

\begin{definition}[Less-than language indicator]
\lean{BoolCircuit.AdjExample.orc}
\label{BoolCircuit.AdjExample.orc}
\leanok
\uses{Complexity.ltLang, Turing.MultiTapeTM.indicator}
Define a Boolean-valued function on a singleton index set and Boolean lists by assigning
to each input list the indicator of its membership in the less-than language.
\end{definition}

\begin{lemma}[Oracle test for a strict inequality]
\lean{BoolCircuit.AdjExample.orc_eq}
\label{BoolCircuit.AdjExample.orc_eq}
\leanok
\uses{BoolCircuit.AdjExample.orc, Complexity.LogProg.bits_injective, Complexity.LogProg.indicator_eq_decide, Turing.pairEncode, Turing.pairEncode_replicate_inj}
\difficulty{4}
For all natural numbers $n$ and $v$, the oracle returns true on the encoding $\langle
1^n,\operatorname{bits}(v)\rangle$ exactly when $v<n$.
\end{lemma}

\begin{lemma}[Updating a constant function]
\lean{BoolCircuit.AdjExample.upd_const}
\label{BoolCircuit.AdjExample.upd_const}
\leanok
\difficulty{3}
For any natural numbers $c$ and $d$, changing the value of the constant function with
value $c$ at the unique element of a one-element domain to $d$ yields the constant
function with value $d$.
\end{lemma}

\begin{lemma}[Bounded-space execution of the comparison program]
\lean{BoolCircuit.AdjExample.loop}
\label{BoolCircuit.AdjExample.loop}
\leanok
\uses{BoolCircuit.AdjExample.Lb, BoolCircuit.AdjExample.leA, BoolCircuit.AdjExample.orc, BoolCircuit.AdjExample.orc_eq, BoolCircuit.AdjExample.upd_const, Complexity.LogProg.AHalt, Complexity.LogProg.AHalt.ret, Complexity.LogProg.AHalt.step, Complexity.LogProg.PreS, Complexity.LogProg.ValidPlain, Complexity.LogProg.astep, Complexity.LogProg.astep_call₁, Complexity.LogProg.bits_injective, Complexity.LogProg.canon_bits, Complexity.LogProg.dbl_replicate, Complexity.LogProg.plainWord, Complexity.LogProg.plainWord_pairEncode, Turing.pairEncode, Turing.pairEncode_eq_dbl}
\difficulty{6}
For any Boolean values $a,b$ and natural numbers $n,i$, for every pair of natural
numbers $c,d$ with $c+d=n$ and $c\le i$, the comparison program started in its test
state with its sole register equal to $c$ halts on input $\langle
1^n,\operatorname{bits}(i)\rangle$, returns the Boolean value of $(i<n\text{ and
}a=\mathrm{true})$ or $(i=n\text{ and }b=\mathrm{true})$, and throughout the run
satisfies the stated program invariant and keeps every register at most $(|x|+1)$, where
$x=\langle 1^n,\operatorname{bits}(i)\rangle$.
\end{lemma}

\begin{theorem}[Logspace membership of the comparison language]
\lean{BoolCircuit.AdjExample.leLang_mem}
\label{BoolCircuit.AdjExample.leLang_mem}
\leanok
\uses{BoolCircuit.AdjExample.leA, BoolCircuit.AdjExample.leLang, BoolCircuit.AdjExample.loop, BoolCircuit.AdjExample.orc, Complexity.LOGSPACE, Complexity.LogProg.AHalt.step, Complexity.LogProg.PreS, Complexity.LogProg.ValidPlain, Complexity.LogProg.arm_decides_poly, Complexity.LogProg.arun, Complexity.LogProg.astep, Complexity.LogProg.bitsVal, Complexity.LogProg.bits_injective, Complexity.LogProg.canon_bits, Complexity.LogProg.canon_eq_bits, Complexity.LogProg.indicator_eq_decide, Complexity.ltLang, Complexity.ltLang_mem, Turing.MultiTapeTM.indicator, Turing.pairEncode, Turing.pairEncode_replicate_inj}
\difficulty{5}
For any two Boolean values $a$ and $b$, the language of bit strings of the form $\langle
1^n,\operatorname{bits}(i)\rangle$ for which either $i<n$ and $a$ is true, or $i=n$ and
$b$ is true, is decidable in logarithmic space. Here $\langle\cdot,\cdot\rangle$ is the
self-delimiting pair encoding.
\end{theorem}

\begin{definition}[Constant-false circuit family]
\lean{BoolCircuit.AdjExample.family}
\label{BoolCircuit.AdjExample.family}
\leanok
\uses{BoolCircuit.DAGCircuitFamily, BoolCircuit.constCircuit}
Define a circuit family that, for every input length $n$, assigns the constant-false
circuit on $n$ inputs.
\end{definition}

\begin{lemma}[Canonical form of the example family]
\lean{BoolCircuit.AdjExample.family_isCanonical}
\label{BoolCircuit.AdjExample.family_isCanonical}
\leanok
\uses{BoolCircuit.AdjExample.family, BoolCircuit.DAGCircuit.size, BoolCircuit.DAGCircuitFamily.IsCanonical, BoolCircuit.constCircuit}
\difficulty{3}
For every input length, the circuit assigned by the example family has its argument
indices in each gate listed in strictly increasing order, and its output vertex index is
one less than the total number of input and gate vertices.
\end{lemma}

\begin{lemma}[Constant output of the circuit family]
\lean{BoolCircuit.AdjExample.family_eval}
\label{BoolCircuit.AdjExample.family_eval}
\leanok
\uses{BoolCircuit.AdjExample.family, BoolCircuit.DAGCircuit.eval}
\difficulty{1}
For every natural number $n$ and every Boolean input $x : \bbr^n \to
\{\mathrm{false},\mathrm{true}\}$, the circuit assigned by the family to input length
$n$ evaluates to false on $x$.
\end{lemma}

\begin{theorem}[Logspace uniformity of the example family]
\lean{BoolCircuit.AdjExample.family_isLogspaceUniform}
\label{BoolCircuit.AdjExample.family_isLogspaceUniform}
\leanok
\uses{BoolCircuit.AdjExample.family, BoolCircuit.AdjExample.family_isCanonical, BoolCircuit.AdjExample.leLang, BoolCircuit.AdjExample.leLang_mem, BoolCircuit.DAGCircuit.size, BoolCircuit.DAGCircuit.vtype, BoolCircuit.DAGCircuitFamily.IsLogspaceUniform, BoolCircuit.DAGCircuitFamily.edgeLang, BoolCircuit.DAGCircuitFamily.isLogspaceUniform_iff, BoolCircuit.DAGCircuitFamily.sizeLang, BoolCircuit.DAGCircuitFamily.typeLang, BoolCircuit.constCircuit, BoolCircuit.constGate}
\difficulty{5}
The example circuit family is logspace-uniform: there is a function mapping binary
strings to binary strings whose output length is bounded by a polynomial in the input
length, whose bit values and output-length predicate are each decidable in logarithmic
space when given an input string and an index, and which, on the unary string of $n$
ones, outputs the encoding of the circuit for length $n$.
\end{theorem}
